Our proof is by decreasing induction on $k$. The initial case, $k=n-1$, is the main result of~\cite{maeno.naito.sagaki:QKideal}, rewritten in determinantal form in Theorem~\ref{thm:toda} below. For the induction step, we use a~result of Kato~\cite{kato2019quantum} which states that there is a
$\K_T(\pt)$-{\em algebra} homomorphism
\begin{gather}
\QK_T(\Fl(n)) \to \QK_T(\Fl(r_1, \dots , r_k;n)); \nonumber\\
\qquad{}\cO^w \mapsto \pi_*(\cO^w) ,
\qquad Q_r \mapsto \begin{cases} 1, & r \notin \{ r_1, \dots, r_k \}, \\
Q_i, & r=r_i ,\end{cases}\label{E:kato-push}
\end{gather}
which extends the usual projection map $\pi_*\colon \K_T(\Fl(n)) \to \K_T(\Fl(\bfr,n))$.
Note that the classical~$\pi_*$ is {\em not} a ring map. (Kato's result is for general complex, simple groups $G$.) We use this to show that the ideal $J_Q$ is contained in the ideal of relations.
For the specialization ${Q_i \mapsto 1}$ to be well defined, one needs to work with
polynomials in $Q_1, \dots, Q_{n-1}$; see Section~\ref{sec:prelim}.
Pushing forward the original Toda relations is not possible, due to poles at $Q_i=1$.
We had to rewrite these relations, and additionally
use an extra identity due to Maeno, Naito, and Sagaki
(cf.~Proposition~\ref{prop:MNS} below), in order for the push forward to be performed.
Finally, it follows from~\cite{GMSXZZ:Nakayama} that the ideal $J_Q$ coincides with the ideal of relations.
